\section{Reconstruction explicite à partir des retours complets}
\label{ret:seccion}

\subsection{Domaine des états et extension bilatérale}

On considère un système de coordonnées du domaine survivant engendré par
APP--TRIT--TPK. Un état de ce système de coordonnées conserve les feuillets opératoriels,
l’orientation tritique, le résidu, le quotient, la retenue, la route,
la frontière, le registre de mémoire et l’origine du calendrier.
La mise à jour agit sur cet état complet.

Soient \(X_0\) un espace compact de Hausdorff non vide et
\(F:X_0\to X_0\) une mise à jour continue et surjective.
Ce sont les hypothèses du résultat dans le système de coordonnées considéré.
L’extension bilatérale de la mise à jour est
\begin{equation}
 \mathscr X_F=
 \bigl\{(x_k)_{k\in\ZZ}\in X_0^{\ZZ}:
                  F(x_k)=x_{k+1}\text{ pour tout }k\in\ZZ\bigr\},
 \label{ret:extension}
\end{equation}
avec la topologie induite par le produit.

\begin{lemma}[Existence et transport bilatéral]
\label{ret:existencia}
L’espace \(\mathscr X_F\) est compact, séparé et non vide.
Chaque projection coordonnée \(\mathscr X_F\to X_0\) est surjective.
Le décalage
\[
 \sigma((x_k)_{k\in\ZZ})=(x_{k+1})_{k\in\ZZ}
\]
est un homéomorphisme de \(\mathscr X_F\).

\end{lemma}
\begin{proof}
Chaque condition \(F(x_k)=x_{k+1}\) définit un fermé du produit
compact de Hausdorff \(X_0^{\ZZ}\).
Un indice \(k_0\) et un état \(z\in X_0\) étant fixés, toute collection
finie de ces conditions, avec \(x_{k_0}=z\), admet une solution :
on complète vers l’avant au moyen de \(F\) et vers l’arrière au moyen de
préimages, qui existent par surjectivité. Les conditions possèdent
donc la propriété d’intersection finie. La compacité produit
une histoire bilatérale avec \(x_{k_0}=z\).

Les décalages \(\sigma\) et
\(\sigma^{-1}((x_k))=(x_{k-1})\) conservent les relations de
\eqref{ret:extension} ; tous deux sont continus par leur action sur les coordonnées.

\end{proof}

\subsection{Échantillonnage tritique et reconstruction de tous les pas}

L’origine du calendrier étant fixée, les instants typés de retour sont
\[
 \mathcal T_+
 =\{k\in\ZZ:1+(k\bmod 9)\in\{1,4,7\}\}=3\ZZ.
\]
À chaque instant, l’état complet \(x_k\) est conservé.
Nous définissons l’espace des registres de retour par

\begin{equation}
 \mathscr Y_+
 =\bigl\{(y_j)_{j\in\ZZ}\in X_0^{\ZZ}:
                 F^3(y_j)=y_{j+1}\text{ pour tout }j\in\ZZ\bigr\}.
 \label{ret:registros}
\end{equation}

\begin{theorem}[Récupération à partir des retours complets]
\label{ret:homeomorfismo}
L'application
\[
 r=\operatorname{Ret}_+:\mathscr X_F\longrightarrow\mathscr Y_+,
 \qquad r(x)_j=x_{3j},
\]
est un homéomorphisme. Son inverse est donné explicitement par
\begin{equation}
 \boxed{
  \bigl(r^{-1}(y)\bigr)_{3j+s}=F^s(y_j),
  \qquad j\in\ZZ,\quad s\in\{0,1,2\}.}
 \label{ret:inversa}
\end{equation}
\end{theorem}
\begin{proof}
Pour \(x\in\mathscr X_F\),
\(F^3(x_{3j})=x_{3j+3}\), de sorte que \(r(x)\in\mathscr Y_+\).
Réciproquement, chaque entier \(k\) possède une unique expression
\(k=3j+s\), avec \(j\in\ZZ\) et \(0\le s<3\), également lorsque
\(k<0\). La formule \eqref{ret:inversa} définit donc toutes les
coordonnées d’une suite \(x\).


Si \(s=0\) ou \(s=1\), on a
\(F(x_{3j+s})=F^{s+1}(y_j)=x_{3j+s+1}\).
À la frontière entre deux blocs,
\[
 F(x_{3j+2})=F^3(y_j)=y_{j+1}=x_{3j+3}.
\]
Ainsi, \(x\in\mathscr X_F\).
La lecture de ses coordonnées d’indice \(3j\) renvoie \(y_j\).
Dans le sens inverse, les relations de \eqref{ret:extension} donnent
\(x_{3j+s}=F^s(x_{3j})\), de sorte que la reconstruction de \(r(x)\)
renvoie \(x\) intégralement.

L’application \(r\) est continue puisqu’il s’agit d’une sélection de coordonnées.
Chaque coordonnée de son inverse est la composition d’une projection
coordonnée avec l’une des applications continues \(I,F,F^2\).
L’inverse est, par conséquent, continu.
\end{proof}

Une coordonnée discrète supplémentaire, par exemple l’indice \(m\in\ZZ\)
d’un système de coordonnées de réalisation, est conservée au moyen de
\[
 \widetilde r:\ZZ\times\mathscr X_F
 \longrightarrow\ZZ\times\mathscr Y_+,
 \qquad \widetilde r(m,x)=(m,r(x)).
\]
Son inverse est \((m,y)\mapsto(m,r^{-1}(y))\).
Le même résultat vaut pour une origine de calendrier fixe \(a\) :
on remplace \(x_{3j}\) et \(x_{3j+s}\) par
\(x_{a+3j}\) et \(x_{a+3j+s}\).

\subsection{Conjugaison du transport et conservation des lectures
}

Soit \(\sigma_+(y)_j=y_{j+1}\) le décalage des registres.
La reconstruction précédente donne

\begin{equation}
 r\sigma^3=\sigma_+r,\qquad
 (r\sigma r^{-1}(y))_j=F(y_j).
 \label{ret:conjugacion}
\end{equation}
L’application de la seconde égalité est réversible dans
\(\mathscr Y_+\) ; son inverse est
\[
 y\longmapsto\bigl(F^2(y_{j-1})\bigr)_{j\in\ZZ}.
\]
En effet, les deux compositions se réduisent à
\(F^3(y_{j-1})=y_j\).
Ainsi, la réversibilité de l’histoire bilatérale est conservée même
lorsque la mise à jour \(F\) sur un état isolé est non injective.
Si le tour nonadique est réalisé comme \(\Gamma_9=\sigma^9\), alors
\begin{equation}
 r\Gamma_9=\sigma_+^3r.
 \label{ret:vuelta-nonadica}
\end{equation}

\begin{corollary}[Lectures conjointes et transport récupéré
]
\label{ret:lecturas}
Soient \(a:\mathscr X_F\to\mathcal R\) et
\(q:\mathscr X_F\to\mathcal E\) deux lecteurs continus du même
état, et soit \(d:\mathscr X_F\to\mathscr X_F\) un homéomorphisme du
domaine calibré. Alors

\[
 a_+=a\circ r^{-1},\qquad
 q_+=q\circ r^{-1},\qquad
 d_+=r\circ d\circ r^{-1}
\]
sont continus, \(d_+\) est un homéomorphisme et
\begin{equation}
 (a,q)=(a_+,q_+)\circ r,\qquad rd=d_+r.
 \label{ret:factorizacion}
\end{equation}
Pour tout entier \(N>0\),
\[
 d^Nx=x\quad\Longleftrightarrow\quad d_+^Nr(x)=r(x).
\]
\end{corollary}
\begin{proof}
On compose les applications avec l’homéomorphisme du
\cref{ret:homeomorfismo}. L’inverse de \(d_+\) est
\(rd^{-1}r^{-1}\). L’identité
\(d_+^Nr=rd^N\), avec l’injectivité de \(r\), prouve la
dernière équivalence.
\end{proof}

L’application à la double lecture prend
\(\mathcal R=R_{\rm HMT}\) et le codomaine incidentiel construit dans
\cref{exc:doble-lectura,exc:limite-inverso}.
Le résultat établit une présentation récupérable des mêmes
histoires, de leurs lecteurs et de leur transport. Son objet est cette
équivalence dynamique de présentations.

